\chapter{Analyse des coûts de transaction}\label{ch:mx:transaction-cost-analysis}

L'ordre d'un gestionnaire de portefeuille a coûté 25 points de base par rapport au prix d'arrivée~; le rapport du courtier indique qu'il a battu le VWAP de trois. Les deux chiffres sont justes, et
seul l'un des deux dit ce que l'ordre a coûté au fonds. L'analyse des coûts de transaction (le livre~7, chapitre~23) est la discipline qui consiste à choisir le chiffre
qui répond à la question posée, de l'expliquer et de comparer des ordres qui n'étaient pas également difficiles à exécuter. Ce chapitre construit un journal d'ordres simulé dans
lequel chaque ordre a été exécuté par deux algorithmes, de sorte que la vraie différence entre eux est connue, puis analyse le journal dont disposerait réellement
un desk, celui où chaque ordre n'a été confronté qu'à un seul algorithme.

\section{Références d'exécution et ce que chacune masque}

\begin{definition}[Référence d'exécution]\label{def:mx:transaction-cost-analysis:benchmark}
Une \emph{référence d'exécution}\index{référence d'exécution} (execution benchmark) est le prix de référence auquel une exécution est mesurée~: le prix de décision ou d'arrivée
(implementation shortfall), le prix moyen pondéré par les volumes sur l'intervalle de l'ordre ou sur la journée, ou le cours de clôture. Le slippage par rapport à elle est le prix
moyen d'exécution moins la référence, signé de sorte qu'un nombre positif soit un coût.
\end{definition}

Le journal simulé couvre 100 jours de \texttt{firm.agentmkt} sur \texttt{firm.exchsim}, chacun d'une durée de 18 minutes, avec un régime de volatilité (la valeur cachée
sautant de 0.1, 0.2 ou 0.4 fois par seconde) et un niveau d'activité (0.7, 1 ou 1.4 fois le flux de base) tiré chaque jour. Chaque jour porte deux \omterm{def:mx:the-almgren-chriss-framework:parent}{ordres parents},
un achat ou une vente de 2\,000 à 12\,000 actions (2~\% à 24~\% du volume sur leur fenêtre), décidés 30 secondes avant leur début, exécutés sur cinq minutes
en dix tranches de 30 secondes, avec une limite à 20 ticks du prix d'arrivée. Deux algorithmes (les calendriers du chapitre~16 sur l'exécuteur du chapitre~17) les exécutent~:
\begin{itemize}
\item \emph{patient}~: tranches égales, chacune placée au meilleur cours affiché puis réajustée, avec une \omterm{def:mx:child-order-placement:risk}{transaction de solde} (clean-up) en fin de tranche~;
\item \emph{urgent}~: un calendrier d'Almgren--Chriss concentré en début de période, avec $\kappa T=3$ (chapitre~14), où chaque tranche est croisée d'un seul coup.
\end{itemize}
Chaque ordre est exécuté par les deux, et la journée est aussi rejouée sans lui, le tout avec le même flux d'ordres exogène (le contrefactuel du chapitre~11). Les prix sont
exprimés en ticks~: un tick vaut un cent sur une action valant 100~USD, soit environ un point de base.

\begin{table}[ht]
\centering
\small
\begin{tabular}{@{}lrrrr@{}}
\toprule
algorithme & vs prix d'arrivée & vs VWAP d'intervalle & vs VWAP de la journée & vs clôture \\
\midrule
patient & 2.42 & $-0.29$ & 1.90 & 2.51 \\
urgent & 3.66 & 1.40 & 3.39 & 4.10 \\
\bottomrule
\end{tabular}
\caption{Slippage moyen par action exécutée (ticks) des 200 ordres simulés sous chaque algorithme, par rapport à quatre \omterm{def:mx:transaction-cost-analysis:benchmark}{références d'exécution}. Données~:
\texttt{mx\_tca.study}.}\label{tab:mx:transaction-cost-analysis:benchmarks}
\end{table}

L'algorithme patient «~bat~» le VWAP d'intervalle de 0.29 ticks tout en coûtant 2.42 par rapport au prix d'arrivée~: le VWAP sur sa propre fenêtre inclut ses
propres transactions et le prix qu'elles ont fait bouger, si bien qu'il évolue avec l'ordre. Par rapport au VWAP de la journée ou à la clôture, les deux algorithmes paraissent meilleurs ou pires selon
ce qu'a fait le reste de la journée. Seul le prix d'arrivée, ou le prix de décision qui le précède, est fixé avant que l'ordre ne touche le marché. Perold (1988)
a opposé le portefeuille théorique, négocié au prix de décision, au portefeuille réel~: l'implementation shortfall.

\section{Estimations des coûts pré-négociation}

\begin{definition}[Estimation des coûts pré-négociation]\label{def:mx:transaction-cost-analysis:pretrade}
Une \emph{estimation des coûts pré-négociation}\index{estimation des coûts pré-négociation} (pre-trade cost estimate) est le coût attendu d'un ordre, calculé avant son envoi à partir de sa taille, du
volume de l'action, de sa volatilité et de son spread, ainsi que d'un modèle d'impact~; elle chiffre la décision de négocier et fixe l'étalon du résultat.
\end{definition}

La loi en racine carrée (le livre~7, chapitre~27, et chapitre~11) donne le modèle~: un demi-spread plus $\eta\,\sigma\sqrt{Q/V}$, avec $\sigma$ la
volatilité du prix sur l'horizon de l'ordre et $Q/V$ sa participation. Le desk estime $\sigma$ pour chaque régime de volatilité à partir de ses journées sans ordres
(7.6, 10.0 et 11.5 ticks sur cinq minutes) et ajuste $\eta$ sur les ordres de l'algorithme patient des 50 premiers jours (avec $-0.5$ tick, le demi-spread
qu'il gagne, en tant qu'ordonnée à l'origine)~: $\eta=1.04$ avec une erreur type de 0.18 (\texttt{firm\_tca.fit\_pretrade}, sur le \texttt{firm.tcost} du livre~7). Sur les
50 jours suivants, le modèle prédit 2.47 ticks par action en moyenne et les ordres coûtent 2.73 (erreur type 0.52)~: juste en moyenne, mais la corrélation
entre la prédiction et le résultat, sur l'ensemble des ordres, n'est que de 0.25 (la \cref{fig:mx:transaction-cost-analysis:calibration}). Le coût d'un ordre isolé est surtout le
mouvement propre du prix~; un modèle pré-négociation est une affirmation sur des moyennes.

\begin{omfigure}
\begin{tikzpicture}
\begin{axis}[width=8.6cm,height=5.8cm,xlabel={\small coût prédit (ticks par action)},ylabel={\small coût réalisé (ticks par action)},
  tick label style={font=\footnotesize},grid=major,grid style={gray!25},xmin=0,xmax=5,ymin=-2,ymax=8,table/col sep=comma]
\addplot[omDef,only marks,mark=*,error bars/.cd,y dir=both,y explicit] table[x=pred,y=real,y error=se]{figdata/microstructure/19-transaction-cost-analysis/calibration.csv};
\addplot[black,dashed,domain=0:5] {x};
\end{axis}
\end{tikzpicture}

\omcaption{Le modèle pré-négociation hors échantillon~: les ordres de l'algorithme patient des 50 derniers jours, répartis en cinq groupes selon le coût prédit, leur coût réalisé moyen
par rapport au prix d'arrivée et son erreur type. Données~: \texttt{mx\_tca.study}.}\label{fig:mx:transaction-cost-analysis:calibration}
\end{omfigure}

\section{Attribution post-négociation}

\begin{definition}[Attribution du slippage]\label{def:mx:transaction-cost-analysis:attribution}
L'\emph{attribution du slippage}\index{attribution du slippage} (slippage attribution) décompose l'implementation shortfall d'un ordre en parties dont la somme le reconstitue~: le délai avant
que l'ordre n'atteigne le marché, le spread payé ou gagné sur chaque exécution, l'impact (le mouvement du prix causé par l'ordre), le timing (le mouvement du prix
que l'ordre n'a pas causé), le coût d'opportunité de ce qui n'a pas été exécuté, et les frais.
\end{definition}

Avec un prix de décision $d$, un prix d'arrivée $a$, des exécutions $q_i$ au prix $p_i$ lorsque le mid valait $m_i$ (juste avant l'exécution) et aurait valu $\tilde m_i$ sans
l'ordre, une cible $X$ et un mid final $e$, le shortfall d'un achat est
\[
\underbrace{\textstyle\sum q_i(a-d)}_{\text{délai}}+\underbrace{\textstyle\sum q_i(p_i-m_i)}_{\text{spread}}+\underbrace{\textstyle\sum q_i(m_i-\tilde m_i-(a-\tilde a))}_{\text{impact}}+\underbrace{\textstyle\sum q_i(\tilde m_i-\tilde a)}_{\text{timing}}+\underbrace{(X-\textstyle\sum q_i)(e-d)}_{\text{opportunité}}+\text{frais},
\]
ce qui correspond au découpage de Perold (\texttt{firm\_markout.shortfall}) en délai, exécution, opportunité et frais, l'exécution étant divisée en trois. La TCA réelle ne voit pas
$\tilde m$~; elle sépare l'impact du timing grâce à un modèle ou au mouvement du marché. Le simulateur le voit, ce qui rend la décomposition ci-dessous exacte.

\omcode{code/firm/tca/firm_tca.py}{44}{60}{Attribution de l'implementation shortfall d'un ordre parent, par action de sa cible~: le mid avant chaque exécution comparé au prix d'exécution (spread), au mid que le même marché aurait eu sans l'ordre (impact), et le mouvement propre de ce mid (timing).}\label{lst:mx:transaction-cost-analysis:attribute}

\begin{omfigure}
\begin{tikzpicture}
\begin{axis}[width=11cm,height=5.6cm,ybar,bar width=8pt,ylabel={\small ticks par action de la cible},
  xtick={0,1,2,3,4,5,6},xticklabels={délai,spread,impact,timing,{opport.},frais,total},
  tick label style={font=\footnotesize},ymin=-1.5,ymax=4.5,grid=major,grid style={gray!25},enlarge x limits=0.08,
  legend columns=2,legend style={font=\footnotesize,at={(0.5,-0.2)},anchor=north,draw=none,fill=none,
  /tikz/every even column/.append style={column sep=8pt}},table/col sep=comma]
\addplot[fill=omDef!70,draw=omDef] table[x=i,y=patient]{figdata/microstructure/19-transaction-cost-analysis/attribution.csv};
\addplot[fill=omProp!70,draw=omProp] table[x=i,y=urgent]{figdata/microstructure/19-transaction-cost-analysis/attribution.csv};
\legend{ordre patient,ordre urgent}
\end{axis}
\end{tikzpicture}

\omcaption{Attribution de l'implementation shortfall des 200 ordres sous chaque algorithme~: les parties s'additionnent exactement au total. Données~:
\texttt{mx\_tca.study}.}\label{fig:mx:transaction-cost-analysis:attribution}
\end{omfigure}

Les deux algorithmes ne se paient pas de la même manière (\cref{fig:mx:transaction-cost-analysis:attribution}). Le patient gagne 0.65 \omterm{def:mx:the-limit-order-book:tick}{pas de cotation} de fourchette et paie 2.41 de
impact~: laisser reposer un gros ordre au meilleur cours affiché pendant cinq minutes maintient le prix de son côté. L'urgent paie 2.43 de fourchette, qui inclut le fait de balayer
le carnet sur ses premières grandes tranches, et 1.13 d'impact. Le timing moyen est proche de zéro (0.59 et $-0.07$), comme il se doit avec des sens aléatoires~; le coût d'opportunité est
faible (99.5~\% et 98.0~\% exécutés dans la limite)~; les frais sont une remise pour l'algorithme patient ($-0.20$) et une commission pour l'urgent (0.29). Au
total, 1.90 contre 3.79 \omterm{def:mx:the-limit-order-book:tick}{pas de cotation} par action, avec des écarts-types de 5.98 et 4.60~: l'algorithme urgent coûte plus cher et varie moins.

\begin{definition}[Retour post-négociation]\label{def:mx:transaction-cost-analysis:reversion}
Le \emph{retour post-négociation}\index{retour post-négociation} (post-trade reversion) est le mouvement du prix après la dernière exécution d'un ordre, signé à l'encontre de celui-ci~: un prix qui
reprend une partie de la poussée de l'ordre (négative) traduit un impact temporaire~; un prix qui continue sa course montre que l'ordre portait une information ou qu'il a provoqué une
\omterm{def:mx:dark-pools-and-blocks:leakage}{fuite d'information} (chapitre~9).
\end{definition}

Deux minutes après la dernière exécution, le prix a repris 3.51 \omterm{def:mx:the-limit-order-book:tick}{pas de cotation} pour l'algorithme patient et 1.06 pour l'urgent (1.37 et 0.07 après 30
secondes). Les dernières exécutions de l'algorithme patient sont souvent ses \omterm{def:mx:child-order-placement:risk}{transactions de solde}, qui poussent le prix en fin de fenêtre~; les dernières
tranches de l'algorithme urgent sont petites. Aucun ordre simulé ne porte d'information, donc aucun prix ne poursuit sa course~; dans les journaux réels, une poursuite après l'achèvement de l'ordre est le signal d'une \omterm{def:mx:dark-pools-and-blocks:leakage}{fuite d'information}.

\section{Comparer entre ordres, courtiers et pairs}

Quel algorithme est le moins cher~? Dans le journal complet la réponse est connue~: en faisant passer chaque ordre par les deux, l'algorithme urgent coûte 1.89 \omterm{def:mx:the-limit-order-book:tick}{pas de cotation} par action de plus,
avec une erreur type de 0.25 regroupée par jour. Le journal du desk est différent. Son routage confiait à l'algorithme urgent les grands ordres des jours volatils plus
souvent (53.5~\% des ordres lui étaient confiés, avec une moyenne de 7\,039 actions contre 4\,126 pour le patient), donc la différence brute dans son journal est de 3.36 \omterm{def:mx:the-limit-order-book:tick}{pas de cotation}
(erreur type de 0.76)~: les ordres difficiles faisaient paraître l'algorithme urgent pire qu'il ne l'est.

\begin{definition}[Comparaison avec un univers de pairs]\label{def:mx:transaction-cost-analysis:peer}
Une \emph{comparaison avec un univers de pairs}\index{comparaison avec un univers de pairs} (peer-universe comparison) mesure le coût d'une exécution par rapport à celui d'ordres comparables passés par d'autres traders,
courtiers ou fonds, après correction de la difficulté de chaque ordre (taille rapportée au volume, volatilité, fourchette, \omterm{def:mx:the-almgren-chriss-framework:urgency}{urgence}), généralement sous la forme du résidu d'un modèle de
coût estimé sur l'univers.
\end{definition}

Ajuster revient à régresser le coût sur l'algorithme et sur la difficulté de l'ordre (la volatilité du régime, la participation, le terme en racine carrée
et son interaction avec l'algorithme), avec des erreurs types regroupées par jour car les deux ordres d'un même jour partagent le marché de ce jour (les erreurs regroupées sont traitées dans le
livre~4, chapitre~16)~; \texttt{firm\_tca.adjusted\_difference} renvoie à la fois le coefficient brut et le coefficient ajusté.

\omcode{code/microstructure/19-transaction-cost-analysis/python/mx_tca.py}{148}{152}{Le routage du desk dans le journal simulé~: la probabilité d'envoyer un ordre à l'algorithme urgent augmente avec sa taille et les jours volatils, ce qui rend la comparaison brute biaisée.}\label{lst:mx:transaction-cost-analysis:assigned}

\begin{omfigure}
\begin{tikzpicture}
\begin{axis}[width=9cm,height=4.4cm,xlabel={\small urgent moins patient (pas de cotation par action)},xmin=-0.5,xmax=5.5,ymin=0.4,ymax=3.6,
  ytick={1,2,3},yticklabels={{vérité (les deux algorithmes)},{ajusté},{brut}},tick label style={font=\footnotesize},grid=major,grid style={gray!25},
  table/col sep=comma]
\addplot[omDef,only marks,mark=*,mark size=2.4pt,error bars/.cd,x dir=both,x explicit]
  table[x=est,y=y,x error plus expr=\thisrow{hi}-\thisrow{est},x error minus expr=\thisrow{est}-\thisrow{lo}]{figdata/microstructure/19-transaction-cost-analysis/estimates.csv};
\addplot[black,dashed] coordinates {(0,0.4) (0,3.6)};
\end{axis}
\end{tikzpicture}

\omcaption{La différence de coût entre les deux algorithmes avec des intervalles à 95~\%~: brute sur le journal du desk, ajustée pour la difficulté sur le même journal, et la
vérité obtenue en faisant passer chaque ordre par les deux. Données~: \texttt{mx\_tca.study}.}\label{fig:mx:transaction-cost-analysis:estimates}
\end{omfigure}

La différence ajustée est de 2.00 \omterm{def:mx:the-limit-order-book:tick}{pas de cotation}, avec un intervalle à 95~\% allant de 0.32 à 3.68 (\cref{fig:mx:transaction-cost-analysis:estimates})~: l'ajustement
supprime l'essentiel du biais (3.36 brut contre 1.89 réel) et l'intervalle couvre la vérité, mais il est large, car 200 ordres, dont l'écart type du coût
est de cinq ou six \omterm{def:mx:the-limit-order-book:tick}{pas de cotation}, ne permettent pas de cerner une différence de deux \omterm{def:mx:the-limit-order-book:tick}{pas de cotation}. Mesurés par rapport au modèle pré-négociation, les ordres de l'algorithme patient
coûtent 0.77 \omterm{def:mx:the-limit-order-book:tick}{pas de cotation} de moins que prévu (erreur type de 0.56) et ceux de l'algorithme urgent 1.50 de plus (0.46)~: la même conclusion, avec l'erreur propre du modèle en plus.
Anand, Irvine, Puckett et Venkataraman (2012) ont établi, sur des données institutionnelles de transactions, que les desks de trading conservent leur performance relative
d'une période à l'autre~; Frazzini, Israel et Moskowitz (2018), sur les exécutions réelles d'un grand gérant, trouvent que les coûts réels de transaction sont inférieurs
d'un ordre de grandeur à ceux qu'annonçaient les études antérieures. Les deux travaux exigeaient une correction de la difficulté de ce type avant toute comparaison.

\begin{dated}{2026-09}{Rapports de meilleure exécution en Europe}\label{dat:mx:transaction-cost-analysis:rts}
MiFID II imposait aux plateformes de négociation de publier des rapports sur la qualité d'exécution («~RTS~27~») et aux entreprises d'investissement de publier chaque année leurs cinq principales plateformes
(«~RTS~28~»). La Directive (EU) 2021/338 a suspendu les rapports périodiques des plateformes jusqu'au 28 février 2023. La Directive (EU) 2024/790 du 28 février 2024, constatant
que les rapports étaient rarement lus et ne permettaient pas de comparaisons pertinentes, a supprimé le rapport sur les cinq principales plateformes et remplacé le rapport sur la qualité d'exécution
par une obligation d'informer chaque client du lieu d'exécution de son ordre~; les États membres disposaient jusqu'au 29 septembre 2025 pour la transposer. Au Royaume-Uni, la
FCA a supprimé les deux rapports à compter du 1 décembre 2021 (policy statement PS21/20).
\end{dated}

\section{Boucler la boucle}

La TCA se rentabilise lorsqu'elle modifie ce que fait le desk. La boucle ici ajuste le coût sur l'algorithme et son interaction avec la difficulté sur les
50 premiers jours du journal du desk, et route chaque ordre des 50 derniers jours vers l'algorithme ayant la prévision la plus basse. Le modèle envoie chaque ordre à
l'algorithme patient~; évalué avec la vérité, cela coûte 2.17 \omterm{def:mx:the-limit-order-book:tick}{pas de cotation} par action sur ces jours, contre 3.24 pour le routage du desk et 3.88 pour l'envoi
de tout confier à l'algorithme urgent. La dispersion plus faible de l'algorithme urgent est la raison pour laquelle un desk pourrait encore le choisir~: un gérant qui pénalise
assez fortement la variance du coût (exercice~6) paie volontiers les 1.9 \omterm{def:mx:the-limit-order-book:tick}{pas de cotation}. La comparaison des algorithmes par expérience plutôt que par régression, avec des ordres
assignés au hasard, est l'algo wheel du chapitre~20~; les raffinements des tests A/B (livre~7, chapitre~21, y compris la réduction de variance par CUPED)
s'appliquent directement.

\section{Tutoriel~: analyser un journal d'ordres}\label{tut:mx:transaction-cost-analysis:tut}

\begin{tutorial}
\textbf{Objectif.} Générer un journal des \omterm{def:mx:the-almgren-chriss-framework:parent}{ordres parents} sur le marché simulé, lancer la TCA contre chaque \omterm{def:mx:transaction-cost-analysis:benchmark}{référence d'exécution}, attribuer le shortfall, estimer et tester un
modèle pré-négociation, et comparer deux algorithmes avec et sans ajustement de la difficulté par rapport à la vérité.
\textbf{État final~:} le \cref{tab:mx:transaction-cost-analysis:benchmarks}, les \cref{fig:mx:transaction-cost-analysis:calibration,fig:mx:transaction-cost-analysis:attribution,fig:mx:transaction-cost-analysis:estimates}
et les nombres des sections 4 et 5.
\begin{tutsteps}
\item \textbf{Journal.} \texttt{mx\_tca.conditions(d)}, \texttt{orders(d)}, \texttt{day\_runs(d)} (le jour sans ordres, patient, urgent),
\texttt{log()}.
\item \textbf{Mesure.} \texttt{firm\_tca.attribute}, \texttt{benchmarks}, \texttt{reversion}.
\item \textbf{Pré-négociation.} \texttt{fit\_pretrade}, \texttt{pretrade}.
\item \textbf{Comparer.} \texttt{assigned(r)}, \texttt{adjusted\_difference}, \texttt{peer\_compare}~; \texttt{study()}~; tracer avec \texttt{fig\_tca.py}.
\end{tutsteps}
\textbf{Pour aller plus loin.} Donner à certains ordres une information (une dérive du prix dans leur sens) et observer le \omterm{def:mx:transaction-cost-analysis:reversion}{retour post-négociation} se transformer en poursuite~; router au
hasard et constater que l'ajustement devient inutile~; ajouter un troisième algorithme.
\end{tutorial}

\section{Construction~: analyse des coûts de transaction}\label{bld:mx:transaction-cost-analysis:tca}

\begin{build}
\bfield{Rôle} La couche de mesure derrière l'algo wheel du chapitre~20, l'\omterm{def:mx:benchmark-algorithms:algo}{algorithme d'exécution} du chapitre~28 et les rapports du desk.
\bfield{Interface} \texttt{attribute(side, target, decision, arrival, fills, mid, cf\_mid, end, t\_arrival)},
\texttt{benchmarks(side, fills, arrival, window, day, close)}, \texttt{reversion(side, t\_end, mid, horizons)}, \texttt{pretrade(half\_spread, eta, sigma, participation)},
\texttt{fit\_pretrade}, \texttt{cluster\_ols(y, X, clusters)}, \texttt{adjusted\_difference(cost, treat, controls, clusters)},
\texttt{peer\_compare(cost, predicted, groups, clusters)}.
\bfield{Règles} Les coûts sont positifs lorsqu'ils sont payés~; l'attribution par action de la cible (le coût d'opportunité compte), les \omterm{def:mx:transaction-cost-analysis:benchmark}{références d'exécution} par action exécutée~; le mid juste avant chaque
exécution~; erreurs regroupées par l'unité qui partage un marché (le jour).
\bfield{Tests d'acceptation} \texttt{code/firm/tca/tests/}~: l'attribution sur un exemple traité à la main, dont la somme égale le shortfall du livre~7~; les \omterm{def:mx:transaction-cost-analysis:benchmark}{références d'exécution} et le \omterm{def:mx:transaction-cost-analysis:reversion}{retour post-négociation}~;
l'ajustement avec erreurs regroupées retrouve un effet injecté que la différence brute surestime~; la \omterm{def:mx:transaction-cost-analysis:peer}{comparaison avec un univers de pairs}.
\bfield{Pour aller plus loin} Attribution sans contrefactuel (impact issu du modèle pré-négociation)~; univers de pairs par taille et tranche de liquidité.
\end{build}

\begin{omsources}
\item A.~F.~Perold, ``The implementation shortfall: paper versus reality'', \emph{Journal of Portfolio Management} 14(3), 1988.
\item A.~Anand, P.~Irvine, A.~Puckett et K.~Venkataraman, ``Performance of institutional trading desks: an analysis of persistence in trading
costs'', \emph{Review of Financial Studies} 25(2), 2012.
\item R.~Kissell, \emph{The Science of Algorithmic Trading and Portfolio Management}, Academic Press, 2014.
\item A.~Frazzini, R.~Israel et T.~J.~Moskowitz, ``Trading costs'', document de travail, SSRN 3229719, 2018.
\item Directive (EU) 2021/338 et Directive (EU) 2024/790 modifiant la Directive 2014/65/EU (MiFID II)~; FCA, PS21/20, 2021.
\end{omsources}

\section{Exercices}

\begin{exercise}[$\star$]\label{exo:mx:transaction-cost-analysis:1}
Un achat de 1\,000 actions est décidé à 100.00 et arrive à 100.02~; 800 actions sont exécutées à un prix moyen de 100.05, le cours termine à 100.10, et les frais s'élèvent à 0.3
centime par action. Calculer les coûts de délai, d'exécution, d'opportunité et de frais, ainsi que le shortfall par action de la cible.
\end{exercise}

\begin{exercise}[$\star$]\label{exo:mx:transaction-cost-analysis:2}
Le VWAP d'intervalle du même ordre valait 100.06. Que dit le rapport VWAP, et que dit le rapport d'arrivée~?
\end{exercise}

\begin{exercise}[$\star$]\label{exo:mx:transaction-cost-analysis:3}
Avec le modèle ajusté ($-0.5$ tick de spread, $\eta=1.04$), quelle est l'estimation pré-négociation pour un ordre à une participation de 10~\% dans le régime intermédiaire
($\sigma=10.0$ ticks)~?
\end{exercise}

\begin{exercise}[$\star\star$]\label{exo:mx:transaction-cost-analysis:4}
Pourquoi l'écart brut entre les algorithmes du journal du desk est-il plus grand que la vérité~?
\end{exercise}

\begin{exercise}[$\star\star$]\label{exo:mx:transaction-cost-analysis:5}
Pourquoi le \omterm{def:mx:transaction-cost-analysis:reversion}{retour post-négociation} de l'algorithme patient est-il plus grand que celui de l'algorithme urgent, alors qu'il négocie plus doucement~?
\end{exercise}

\begin{exercise}[$\star\star$]\label{exo:mx:transaction-cost-analysis:6}
Un gestionnaire minimise le coût espéré plus $\lambda$ fois sa variance (par action, en ticks). Au-delà de quelle valeur de $\lambda$ l'écart type plus faible de l'algorithme urgent
justifie-t-il sa moyenne plus élevée~?
\end{exercise}

\begin{exercise}[$\star\star\star$]\label{exo:mx:transaction-cost-analysis:7}
\emph{Programmation.} Ajuster la comparaison sur la seule participation. Quelle différence et quel intervalle obtient-on, et pourquoi~?
\end{exercise}

\begin{exercise}[$\star\star\star$]\label{exo:mx:transaction-cost-analysis:8}
\emph{Trouver la faille.} «~L'algorithme patient a battu le VWAP de 0.3 tick en moyenne, donc ses ordres n'ont rien coûté à nos clients.~»
\end{exercise}

\section{Problème~: mauvais courtier ou journée difficile~?}

\begin{problem}[{Problème du week-end --- mauvais courtier ou journée difficile~?}]\label{pb:mx:transaction-cost-analysis:1}
Le rapport TCA du desk indique que l'algorithme urgent coûte 3.4 ticks de plus par action que l'algorithme patient. Le courtier qui le fournit affirme que ses ordres étaient plus difficiles à exécuter.

\medskip\noindent\textbf{Partie I --- Références.}
\begin{enumerate}
\item Définir les quatre références du chapitre.
\item Donner le slippage de chaque algorithme par rapport à celles-ci.
\item Pourquoi l'algorithme patient bat-il son VWAP d'intervalle~?
\item Quelle référence répond à la question «~combien l'ordre a-t-il coûté au fonds~?~», et pourquoi~?
\end{enumerate}

\medskip\noindent\textbf{Partie II --- Estimations et attribution.}
\begin{enumerate}[resume]
\item Énoncer le modèle pré-négociation et son coefficient ajusté.
\item Quelle est sa qualité de prévision hors échantillon, en moyenne et ordre par ordre~?
\item Écrire l'attribution et montrer qu'elle se somme à l'implementation shortfall.
\item Donner l'attribution de chaque algorithme.
\item Que dit le retour après la dernière exécution~?
\end{enumerate}

\medskip\noindent\textbf{Partie III --- La comparaison.}
\begin{enumerate}[resume]
\item Décrire le routage du desk et les ordres reçus par chaque algorithme.
\item Donner l'écart brut et son erreur type.
\item Écrire la régression ajustée et expliquer le regroupement des observations.
\item \emph{Énoncer le résultat nommé}~: l'écart de coût ajusté de la difficulté entre les deux algorithmes, avec son intervalle de confiance, à comparer à
l'écart brut.
\item Comment se situe-t-il par rapport à la vérité, et pourquoi un desk réel ne peut-il pas voir la vérité~?
\item Que donne la comparaison avec le modèle pré-négociation~?
\end{enumerate}

\medskip\noindent\textbf{Partie IV --- Boucler la boucle.}
\begin{enumerate}[resume]
\item Quel routage le modèle ajusté recommande-t-il, et combien permet-il d'économiser~?
\item Pourquoi un gestionnaire pourrait-il tout de même choisir l'algorithme urgent~?
\item Comment un routage randomisé modifierait-il l'analyse~?
\item Mauvais courtier ou journée difficile~: que répondez-vous~?
\item En une phrase~: qu'est-ce qui rend deux coûts d'exécution comparables~?
\end{enumerate}
\end{problem}

\section{Questions d'entretien}

\begin{interviewq}[$\star$ \iqroles{trader}]\label{iq:mx:transaction-cost-analysis:1}
Votre ordre a battu le VWAP mais a coûté 20 points de base par rapport à l'arrivée. Expliquez les deux chiffres au gérant de portefeuille.
\end{interviewq}

\begin{interviewq}[$\star\star$ \iqroles{chercheur}]\label{iq:mx:transaction-cost-analysis:2}
Comment répartiriez-vous le slippage d'un ordre entre spread, impact et délai, sans contrefactuel~?
\end{interviewq}

\begin{interviewq}[$\star\star$ \iqroles{chercheur}]\label{iq:mx:transaction-cost-analysis:3}
Les coûts moyens de deux courtiers diffèrent de 3 points de base. De quoi avez-vous besoin avant de conclure que l'un est meilleur~?
\end{interviewq}

\begin{interviewq}[$\star\star$ \iqroles{risque}]\label{iq:mx:transaction-cost-analysis:4}
Que vous apprend le \omterm{def:mx:transaction-cost-analysis:reversion}{retour post-négociation}, et comment l'utiliseriez-vous pour détecter une \omterm{def:mx:dark-pools-and-blocks:leakage}{fuite d'information}~?
\end{interviewq}

\begin{interviewq}[$\star\star$ \iqroles{développeur}]\label{iq:mx:transaction-cost-analysis:5}
Concevez les données qu'un système de TCA doit enregistrer pour chaque \omterm{def:mx:the-almgren-chriss-framework:parent}{ordre parent} et chaque \omterm{def:mx:the-almgren-chriss-framework:parent}{ordre enfant}.
\end{interviewq}

\begin{interviewq}[$\star\star\star$ \iqroles{banque}]\label{iq:mx:transaction-cost-analysis:6}
Le prestataire de TCA d'un client classe vos algorithmes en dernier dans son univers de pairs. Comment examinez-vous cette affirmation~?
\end{interviewq}
